\chapterbackmatter{Weinberg Exercises}

\begin{enumerate}
\WeinbergExercise{1} Consider the theory of a real scalar field $\phi$, with interaction
(in the interaction picture)
$H_I=g\int \dd^3x\,\phi(x)^3/3!$.
Calculate the connected $S$-matrix element for scalar--scalar scattering to
second order in $g$, doing all integrals. Use the results to calculate the
differential cross-section for scalar--scalar scattering in the
center-of-mass system.

\WeinbergExercise{2} Consider a theory involving a neutral scalar field $\phi(x)$ for a
boson $B$ and a complex Dirac field $\psi(x)$ for a fermion $F$, with
interaction (in the interaction picture)
$H_I=ig\int \dd^3x\,\bar\psi(x)\gamma_5\psi(x)\phi(x)$.
Draw all the connected order-$g^2$ Feynman diagrams and calculate the
corresponding $S$-matrix elements for the processes
$F^c+B\to F^c+B$, $F+F^c\to F+F^c$, and $F^c+F\to B+B$
(where $F^c$ is the antiparticle of $F$). Do all integrals.

\WeinbergExercise{3} Consider the theory of a real scalar field $\phi(x)$, with interaction
$H_I=g\int \dd^3x\,\phi(x)^4/4!$.
Calculate the $S$-matrix for scalar--scalar scattering to order $g$, and use
the result to calculate the differential scattering cross-section. Calculate
the correction terms in the $S$-matrix for scalar--scalar scattering to order
$g^2$, expressing the result as an integral over a single four-momentum, but
do all $x$-integrals.

\WeinbergExercise{4} What is the contribution in Feynman diagrams from the contraction of
the derivative $\partial_\mu\psi_\ell(x)$ of a Dirac field with the adjoint
$\bar\psi_m(y)$ of the field?

\WeinbergExercise{5} Use the theorem of Section 6.4 to give expressions for the vacuum
expectation values of Heisenberg picture operators
$\bra{\Omega}\Phi(x)\ket{\Omega}$ and
$\bra{\Omega}T\{\Phi(x),\Phi(y)\}\ket{\Omega}$ in the theory of Problem 1,
to orders $g$ and $g^2$, respectively.
\end{enumerate}
